\documentclass[12pt]{article}
\usepackage{geometry}
\usepackage{fontspec}
\usepackage{amsmath,amssymb}
\usepackage{unicode-math}
\usepackage{xcolor}
\usepackage{booktabs}

\geometry{a4paper, top=2.5cm, bottom=2.5cm, left=2.5cm, right=2.5cm}

\definecolor{charcoal}{HTML}{36454F}
\setmainfont{texgyrepagella}[
  Extension = .otf,
  UprightFont = *-regular,
  BoldFont = *-bold,
  ItalicFont = *-italic,
  BoldItalicFont = *-bolditalic,
  Color=charcoal,
]
\setmathfont{texgyrepagella-math.otf}

\setlength{\parindent}{0pt}
\setlength{\parskip}{6pt}

\title{Entropy and Cross-Entropy at Inference Time}
\author{}
\date{}

\begin{document}
\maketitle
\vspace{-2.5em}

\section*{Setting}

A model with parameters $\theta$ takes an input $x$ and gives a
probability distribution $q_\theta(k \mid x)$ over the classes
$k = 1, \dots, K$. In this note, all logarithms are natural logarithms,
so all quantities are measured in nats. If we use $\log_2$ instead, we
get bits. For example, a fair coin has entropy
$\ln 2 \approx 0.693$ nats, which is exactly 1 bit.

In the cross-entropy formula, $p$ is the \emph{target} distribution and
$q$ is the model's \emph{predicted} distribution. The target is what we
compare against; the prediction is what the model gives us. We will see
below what the target looks like for a coin and for labelled data.

\section*{What does entropy measure?}

In this section we use $\log_2$, so entropy is measured in bits. At the
end of the section we explain how bits and nats are related.

Let us start with a fair coin, which lands Heads with probability
$0.5$ and Tails with probability $0.5$. Its entropy is
\begin{equation}
  H = -0.5\log_2 0.5 - 0.5\log_2 0.5
    = 0.5 + 0.5 = 1 \text{ bit per toss}.
  \label{eq:fair}
\end{equation}
This matches our common sense. Before a fair toss, both outcomes are
equally likely, so we have no way to guess better than chance. Writing
Heads as $1$ and Tails as $0$, we need exactly one bit per toss, and no
lossless method can do better on average. For a coin, $1$ bit is the
largest possible entropy: a fair coin is as unpredictable as a coin
can be.

Now take a coin that lands Heads with probability $0.6$ and Tails with
probability $0.4$. Its entropy is
\begin{equation}
  H = -0.6\log_2 0.6 - 0.4\log_2 0.4
    \approx 0.971 \text{ bits per toss}.
  \label{eq:biased}
\end{equation}
This does not mean that one toss can be recorded using a fraction of a
bit: recording a single Heads or Tails result requires one bit. Entropy
is an average over many tosses. If we encode a long sequence of
independent tosses using an efficient lossless method, the average
length can approach $0.971$ bits per toss, but it cannot go below that.
The coin is slightly predictable (Heads comes up more often), and a
good encoding takes advantage of this.

\subsection*{Bits and nats measure the same thing}

Bits and nats are two units for the same quantity, just like kilometres
and miles are two units for the same distance. Changing the base of the
logarithm only multiplies every logarithm by the same constant:
\begin{equation}
  \log_2 a = \frac{\ln a}{\ln 2}
  \qquad\text{for every } a > 0.
  \label{eq:changebase}
\end{equation}
Since entropy is a sum of such logarithms, the whole entropy gets
multiplied by the same constant. So
\begin{equation}
  H \text{ in nats} = (H \text{ in bits}) \times \ln 2,
  \qquad
  H \text{ in bits} = \frac{H \text{ in nats}}{\ln 2},
  \label{eq:convert}
\end{equation}
where $\ln 2 \approx 0.693$. Our two coins in both units:

\begin{center}
\begin{tabular}{lcc}
  \toprule
  Coin & Entropy in bits & Entropy in nats \\
  \midrule
  Fair, $(0.5,\ 0.5)$   & $1$             & $\ln 2 \approx 0.693$ \\
  Biased, $(0.6,\ 0.4)$ & $\approx 0.971$ & $\approx 0.673$ \\
  \bottomrule
\end{tabular}
\end{center}

The ordering never changes: whichever unit we use, the biased coin has
less entropy than the fair coin. Only the numbers change.

\textbf{Which unit should we use?} Bits are easier to picture when we
talk about encoding, because a bit is one $0$ or $1$. That is why we
used $\log_2$ in this section. But in deep learning we use $\ln$:
libraries such as PyTorch compute cross-entropy with $\ln$, so the loss
values you see during training are in nats, and the derivative of $\ln z$
is simply $1/z$, which keeps the gradients simple. So from
the next section onwards, we go back to $\ln$ and nats.

\section*{Inference}

In production, $\theta$ is fixed. For every input $x$, the model gives
$q_\theta(\cdot \mid x)$, and the system uses it either to make a
prediction (for example, pick the class with the highest probability)
or to sample an outcome. Ordinary prediction does not need any loss to
be computed, and $\theta$ is not updated.

\section*{Entropy of a prediction}

The output distribution has an entropy:
\begin{equation}
  H\big(q_\theta(\cdot \mid x)\big)
  = -\sum_{k=1}^{K} q_\theta(k \mid x)\,\ln q_\theta(k \mid x).
  \label{eq:entropy}
\end{equation}
To compute this, we need only the model's output. We do not need any
label. It tells us how spread out the prediction is for this particular
input. If all the probability sits on one class, the entropy is zero.
If the prediction is uniform over all $K$ classes, the entropy takes its
largest value, $\ln K$.

So an application can compute \eqref{eq:entropy} for each input and use
it as a confidence signal. For example, predictions with high entropy
can be sent to a human for review. But please note: low entropy only
means the model is confident. It does not mean the model is correct.

\section*{Cross-entropy}

Cross-entropy compares the model's distribution $q$ with a target
distribution $p$:
\begin{equation}
  H(p, q) = -\sum_{k} p(k)\,\ln q(k).
  \label{eq:crossentropy}
\end{equation}

\subsection*{First, the coin}

For a coin, we can write down $p$ fully. Suppose we toss a fair coin
before a cricket match, so $p = (0.5,\ 0.5)$ for (heads, tails). Suppose
our model wrongly says the coin is biased, $q = (0.98,\ 0.02)$. Then
\begin{equation}
  H(p, q) = -0.5\,\ln 0.98 \;-\; 0.5\,\ln 0.02
  \approx 0.010 + 1.956 = 1.966 \text{ nats}.
  \label{eq:coin}
\end{equation}
If the model were correct, that is $q = p$, we would get only
$\ln 2 \approx 0.693$ nats. The extra $1.966 - 0.693 = 1.273$ nats is
the price we pay for using the wrong model. This extra amount is the
KL divergence, $\mathrm{KL}(p \,\|\, q)$.

We can get the same number $1.966$ from data, without writing $p$ down
at all. Toss the coin many times. For each toss, look at the face that
actually came up, take $-\ln$ of the probability the model gave to that
face, and then take the average over all tosses. As the number of
tosses grows, this average comes closer and closer to $H(p, q)$. So
$p$ enters only through what we actually observed.

\subsection*{Now, labelled data}

With labelled data, the situation is a little different. For each input
$x$ we see only one outcome, namely its label $y$. We do not know the
full distribution of labels for that $x$. So, as the target, we use the
\emph{one-hot} distribution on $y$: probability $1$ on class $y$ and
$0$ on every other class. With this target, \eqref{eq:crossentropy}
becomes
\begin{equation}
  \ell(x, y; \theta) = -\ln q_\theta(y \mid x).
  \label{eq:perexample}
\end{equation}
In simple words: the loss for one example is $-\ln$ of the probability
the model gave to the correct label.

Over a labelled set $\{(x_i, y_i)\}_{i=1}^{N}$, the reported loss is
the average
\begin{equation}
  \mathcal{L}(\theta)
  = \frac{1}{N} \sum_{i=1}^{N} -\ln q_\theta(y_i \mid x_i),
  \label{eq:loss}
\end{equation}
which is the average cross-entropy between each one-hot target and the
model's prediction for that input.

Notice that \eqref{eq:loss} is exactly the same recipe as the coin:
look at what actually happened, take $-\ln$ of the probability the
model gave to it, and average. Averaging over many inputs here plays
the same role that averaging over many tosses played for the coin.

\subsection*{Why $p$ and $q$ cannot be swapped}

The order of $p$ and $q$ in \eqref{eq:crossentropy} matters. Suppose
$p$ is one-hot on $y$ and we swap the roles, giving the reversed
expression
\begin{equation}
  -\sum_{k} q(k)\,\ln p(k).
  \label{eq:reversed}
\end{equation}
This is infinite whenever $q$ gives positive probability to any class
other than $y$, because $p(k) = 0$ for those classes and
$\ln 0 = -\infty$.\footnote{Here we use the usual convention $0 \cdot \ln 0 = 0$.
A model that ends with a softmax gives $q(k) > 0$ for every class, so
in practice the reversed expression is always infinite.}

\section*{Training, evaluation, and inference}

\begin{itemize}
  \item \textbf{Training} minimises $\mathcal{L}(\theta)$ in
    \eqref{eq:loss} with respect to $\theta$. This changes $q_\theta$.
  \item \textbf{Evaluation} computes the same $\mathcal{L}(\theta)$ on
    held-out labelled data, keeping $\theta$ fixed. Computing a loss
    does not by itself mean we are training.
  \item \textbf{Inference} uses $q_\theta(\cdot \mid x)$ to predict or
    to sample. At that moment the true labels are usually not
    available, so label-based cross-entropy cannot be computed then.
    But the entropy \eqref{eq:entropy} of each prediction can always be
    computed. If labels come in later, computing $\mathcal{L}(\theta)$
    on them tells us whether the fixed model still fits the new data.
\end{itemize}

\end{document}